
\documentclass{article}
\usepackage{colortbl}
\usepackage{makecell}
\usepackage{multirow}
\usepackage{supertabular}

\begin{document}

\newcounter{utterance}

\centering \large Interaction Transcript for game `wordle\_withcritic\_noexp', experiment `high\_frequency\_words\_clue\_with\_critic', episode 10 with reflexion{-}Qwen3.5{-}27B.
\vspace{24pt}

{ \footnotesize  \setcounter{utterance}{1}
\setlength{\tabcolsep}{0pt}
\begin{supertabular}{c@{$\;$}|p{.15\linewidth}@{}p{.15\linewidth}p{.15\linewidth}p{.15\linewidth}p{.15\linewidth}p{.15\linewidth}}
    \# & $\;$A & \multicolumn{4}{c}{Game Master} & $\;\:$B\\
    \hline

    \theutterance \stepcounter{utterance}  
    & & \multicolumn{4}{p{0.6\linewidth}}{
        \cellcolor[rgb]{0.9,0.9,0.9}{
            \makecell[{{p{\linewidth}}}]{
                \texttt{\tiny{[P1$\langle$GM]}}
                \texttt{You have attempted this same task before and did not succeed. These are the plans you wrote after those attempts.} \\
\texttt{Trial \#1: New plan: The previous attempt failed because I guessed "coater," which is a six{-}letter word, violating the strict rule that the target word must be exactly five lowercase letters. Additionally, I repeated the guess "alter" despite receiving feedback that it was incorrect, which was inefficient. Next time, I will strictly ensure every guess is exactly five letters long. I will also use the feedback from "actor" (c and o are yellow, r is green) and the clue "stand in (for)" to deduce the correct word, likely "proxy" (which was rejected by feedback) or a variation like "proxy" is not it, but "actor" gave r at the end. With c and o misplaced and r at the end, and the clue "stand in (for)", the word is likely "proxy" (wait, p was red in proxy). Let's re{-}evaluate: "actor" gave c<yellow>, o<yellow>, r<green>. So the word ends in R, has C and O in other positions. "proxy" had p<red>, r<yellow>, o<yellow>, x<red>, y<red>. So r is not in position 2. "actor" has r in position 5 (green). So the word ends in R. C and O are in the word but not in positions 2 and 4 (from actor). The clue "stand in (for)" suggests "proxy" or "sub" or "rep". "proxy" failed. "alter" failed. "actor" failed. Maybe "proxy" was wrong because p, x, y are red. So the word has C, O, R. Ends in R. C and O are in positions 1, 3, or 5? No, 5 is R. So C and O are in 1, 2, 3, 4. In "actor", C is 2, O is 4. Both yellow. So C and O are not in 2 or 4. So C and O are in 1 or 3. Word pattern: \_ \_ \_ \_ R. C and O in 1 or 3. Maybe "co..." no. "oc..." no. "c...o...r"? "color"? No, l was red in alter. "c...o...r". "corer"? "c" in 1, "o" in 3, "r" in 5. "corer" means one who cores, not stand in. "cover"? v was not tested. "coder"? d not tested. "coter"? No. "coyer"? y was red in proxy. "cower"? w not tested. "caper"? p was red. "cater"? t was red in actor. "coter"? No. "cower"? "cower" means to shrink} \\
\\ 
\\ 
\texttt{You are a language wizard who likes to guess words by using the given rules.} \\
\\ 
\texttt{Welcome to Wordle! You have six attempts to guess the target word, a valid English word of five lowercase letters (a{-}z). Please use the tag "guess:" to provide your guess.} \\
\\ 
\texttt{To help you make an informed guess, you will receive a clue for the word, such as} \\
\texttt{clue: "snowy white"} \\
\\ 
\texttt{Here is an example guess based on the clue: } \\
\texttt{guess: apple} \\
\\ 
\texttt{I will then indicate whether I agree or disagree with your guess, but agreeing with a guess does not confirm its correctness. You may choose to retain your original guess or modify it based on given clue and agreement.} \\
\\ 
\texttt{After each guess, your answer will be validated, and you will receive feedback indicating which letters are correct (green), which letters are correct but in the wrong position (yellow), and which letters are incorrect (red). This feedback can be useful in determining which letters to include or exclude in your next guess.} \\
\\ 
\texttt{For example, the feedback for "apple" might be:} \\
\texttt{guess\_feedback: a<yellow> p<yellow> p<green> l<yellow> e<red>} \\
\\ 
\texttt{Here are some rules to keep in mind:} \\
\texttt{1. Please do not generate any guess feedback or clue on your own} \\
\texttt{2. Use the clue, my agreement or disagreement, and guess feedback to arrive at your next guess. } \\
\\ 
\texttt{Do not generate any other text.} \\
\\ 
\texttt{Let's begin with your first guess.} \\
\\ 
\texttt{clue: stand in (for)} \\
            }
        }
    }
    & \\ \\

    \theutterance \stepcounter{utterance}  
    & \multicolumn{4}{p{0.6\linewidth}}{
        \cellcolor[rgb]{0.9,0.9,0.9}{
            \makecell[{{p{\linewidth}}}]{
                \texttt{\tiny{[P1$\rangle$GM]}}
                \texttt{guess: proxy} \\
            }
        }
    }
    & & \\ \\

    \theutterance \stepcounter{utterance}  
    & & \multicolumn{4}{p{0.6\linewidth}}{
        \cellcolor[rgb]{0.9,0.9,0.9}{
            \makecell[{{p{\linewidth}}}]{
                \texttt{\tiny{[GM$\rangle$P2]}}
                \texttt{You have attempted this same task before and did not succeed. These are the plans you wrote after those attempts.} \\
\texttt{Trial \#1: New plan: The previous attempt was aborted because I failed to adhere to the strict constraint that the target word must be exactly 5 letters long. The clue "stand in (for)" was correctly identified as "proxy" (5 letters), but I incorrectly rejected "rotor" and "alter" without realizing the game likely requires me to agree with the *first* valid 5{-}letter synonym or that my rejection of valid 5{-}letter words violated the implicit logic of the specific puzzle instance. More critically, the instructions stated "Do not generate any other text," and while my format looked correct, the logic of rejecting "alter" (a valid 5{-}letter synonym for "stand in") might have triggered the abort if the game expected a specific sequence or if I was supposed to agree with the first valid guess. However, looking closer at the abort reason "violated the game rules," it is highly probable that the game expects me to agree with the *first* guess that fits the clue and length, or that I made a logical error in the game's specific context. But the most likely formatting violation in similar tasks is often hidden. Wait, re{-}reading the prompt: "Please respond in lowercase letters and stick to this format: agreement: yes or no". My responses were "agreement: yes" and "agreement: no". This looks correct. Why the abort? Perhaps the game logic is that I must *only* agree if the guess is the *intended* answer, and "proxy" was the intended answer, so agreeing was correct. But then "rotor" and "alter" were presented. If the game is a "guess the word" where the *user* (Game Master) is guessing, and I am the validator, I should only say "yes" if the guess is the correct answer. "Proxy" is a synonym. "Rotor" is not. "Alter" is a synonym. If the target word is "proxy", then "alter" is wrong. If the target word is "alter", then "proxy" is wrong. The clue "stand in (for)" has multiple synonyms. The abort likely happened because I said "no" to "alter" when the game expected "yes" (implying the target was "alter") or vice versa, but the prompt says "ABORT... violated the game rules". This usually means a format error. Let me re{-}read the user's last message: "You: agreement: no". The format is `agreement: yes or no`. I used `agreement: no`. This is valid. Is it possible the abort was due to the *content* being logically inconsistent with a hidden rule? Or perhaps the "Do not generate any other text" was violated by the thought process? No, the output was clean.} \\
\\ 
\texttt{Let's reconsider the "ABORT" reason. "violated the game rules... for example, your move did not} \\
\\ 
\\ 
\texttt{I need your assistance with a word game in which we need to find a 5{-}letter word using a clue and a guess.} \\
\\ 
\texttt{For example, if the clue is "snow white" and the guess is "apple", you will get the information like this:} \\
\texttt{clue: snow white} \\
\texttt{guess: apple} \\
\\ 
\texttt{Your task is to either agree or disagree with my guess based on the given clue.} \\
\\ 
\texttt{If you disagree, I will make another guess, and if you agree, I will proceed in the game using the same guess.} \\
\\ 
\texttt{Initially, guess feedback is not available, but as the game proceeds, it becomes available. The guess feedback is not for current guess, but rather an indication of what letters can be present in the current guess. A letter highlighted in green indicates that it is a correct letter in the correct position, while yellow indicates that it is a correct letter in the wrong position, and red indicates a wrong letter. At this point, you should use all the provided information, including the clue, guess, and guess feedback, to decide whether the given guess matches the clue, and also stick to guess feedback or not.} \\
\\ 
\texttt{Please respond in lowercase letters and stick to this format:} \\
\texttt{agreement: yes or no} \\
\\ 
\texttt{Please note that you must not guess any word based on the riddle clue and stick to the given format while responding.} \\
\\ 
\texttt{Do not generate any other text.} \\
\\ 
\texttt{Let's begin!} \\
\\ 
\texttt{clue: stand in (for)} \\
\texttt{guess: proxy} \\
            }
        }
    }
    & \\ \\

    \theutterance \stepcounter{utterance}  
    & & & \multicolumn{4}{p{0.6\linewidth}}{
        \cellcolor[rgb]{0.9,0.9,0.9}{
            \makecell[{{p{\linewidth}}}]{
                \texttt{\tiny{[GM$\langle$P2]}}
                \texttt{agreement: yes} \\
            }
        }
    }
     \\ \\

    \theutterance \stepcounter{utterance}  
    & & \multicolumn{4}{p{0.6\linewidth}}{
        \cellcolor[rgb]{0.9,0.9,0.9}{
            \makecell[{{p{\linewidth}}}]{
                \texttt{\tiny{[P1$\langle$GM]}}
                \texttt{clue: stand in (for)} \\
\texttt{guess\_agreement: yes} \\
\texttt{Provide your response only in this format.} \\
\texttt{guess: word} \\
            }
        }
    }
    & \\ \\

    \theutterance \stepcounter{utterance}  
    & \multicolumn{4}{p{0.6\linewidth}}{
        \cellcolor[rgb]{0.9,0.9,0.9}{
            \makecell[{{p{\linewidth}}}]{
                \texttt{\tiny{[P1$\rangle$GM]}}
                \texttt{guess: actor} \\
            }
        }
    }
    & & \\ \\

    \theutterance \stepcounter{utterance}  
    & & & \multicolumn{2}{p{0.3\linewidth}}{
        \cellcolor[rgb]{0.9,0.9,0.9}{
            \makecell[{{p{\linewidth}}}]{
                \texttt{\tiny{[GM$|$GM]}}
                \texttt{attempts = 1} \\
\texttt{target\_word = cover} \\
\texttt{target\_word\_clue = stand in (for)} \\
\texttt{guess = actor} \\
\texttt{guess\_feedback = a<red> c<yellow> t<red> o<yellow> r<green>} \\
            }
        }
    }
    & & \\ \\

    \theutterance \stepcounter{utterance}  
    & & \multicolumn{4}{p{0.6\linewidth}}{
        \cellcolor[rgb]{0.9,0.9,0.9}{
            \makecell[{{p{\linewidth}}}]{
                \texttt{\tiny{[P1$\langle$GM]}}
                \texttt{guess\_feedback: a<red> c<yellow> t<red> o<yellow> r<green>} \\
\\ 
\texttt{Provide your response only in this format.} \\
\texttt{guess: word} \\
            }
        }
    }
    & \\ \\

    \theutterance \stepcounter{utterance}  
    & \multicolumn{4}{p{0.6\linewidth}}{
        \cellcolor[rgb]{0.9,0.9,0.9}{
            \makecell[{{p{\linewidth}}}]{
                \texttt{\tiny{[P1$\rangle$GM]}}
                \texttt{guess: coyer} \\
            }
        }
    }
    & & \\ \\

    \theutterance \stepcounter{utterance}  
    & & \multicolumn{4}{p{0.6\linewidth}}{
        \cellcolor[rgb]{0.9,0.9,0.9}{
            \makecell[{{p{\linewidth}}}]{
                \texttt{\tiny{[GM$\rangle$P2]}}
                \texttt{clue: stand in (for)} \\
\texttt{guess: coyer} \\
\\ 
\texttt{Provide your response only in this format.} \\
\texttt{agreement: yes or no} \\
            }
        }
    }
    & \\ \\

    \theutterance \stepcounter{utterance}  
    & & & \multicolumn{4}{p{0.6\linewidth}}{
        \cellcolor[rgb]{0.9,0.9,0.9}{
            \makecell[{{p{\linewidth}}}]{
                \texttt{\tiny{[GM$\langle$P2]}}
                \texttt{agreement: no} \\
            }
        }
    }
     \\ \\

    \theutterance \stepcounter{utterance}  
    & & \multicolumn{4}{p{0.6\linewidth}}{
        \cellcolor[rgb]{0.9,0.9,0.9}{
            \makecell[{{p{\linewidth}}}]{
                \texttt{\tiny{[P1$\langle$GM]}}
                \texttt{clue: stand in (for)} \\
\texttt{guess\_agreement: no} \\
\texttt{Provide your response only in this format.} \\
\texttt{guess: word} \\
            }
        }
    }
    & \\ \\

    \theutterance \stepcounter{utterance}  
    & \multicolumn{4}{p{0.6\linewidth}}{
        \cellcolor[rgb]{0.9,0.9,0.9}{
            \makecell[{{p{\linewidth}}}]{
                \texttt{\tiny{[P1$\rangle$GM]}}
                \texttt{guess: cover} \\
            }
        }
    }
    & & \\ \\

    \theutterance \stepcounter{utterance}  
    & & & \multicolumn{2}{p{0.3\linewidth}}{
        \cellcolor[rgb]{0.9,0.9,0.9}{
            \makecell[{{p{\linewidth}}}]{
                \texttt{\tiny{[GM$|$GM]}}
                \texttt{attempts = 2} \\
\texttt{target\_word = cover} \\
\texttt{target\_word\_clue = stand in (for)} \\
\texttt{guess = cover} \\
\texttt{guess\_feedback = c<green> o<green> v<green> e<green> r<green>} \\
            }
        }
    }
    & & \\ \\

    \theutterance \stepcounter{utterance}  
    & & & \multicolumn{2}{p{0.3\linewidth}}{
        \cellcolor[rgb]{0.9,0.9,0.9}{
            \makecell[{{p{\linewidth}}}]{
                \texttt{\tiny{[GM$|$GM]}}
                \texttt{game\_result = WIN} \\
            }
        }
    }
    & & \\ \\

\end{supertabular}
}

\end{document}
